%% Allures des courbes du MOUVEMENT UNIFORMÉMENT FREINÉ (décéléré) :
%% position qui s'aplatit, vitesse affine décroissante, accélération constante négative.
\documentclass[tikz,border=4pt]{standalone}
\usepackage{liaisons}
\begin{document}
\begin{tikzpicture}[x=1cm,y=1cm, scale=0.92,
  axe/.style={-{Stealth[length=5pt]}, line width=0.6pt},
  courbe/.style={solideA, line width=1.3pt, line cap=round}]

%% \petitrepere{décalage}{grandeur en ordonnée}{bas de l'axe}{tracé}
\newcommand\petitrepere[4]{%
\begin{scope}[shift={(#1,0)}]
  \draw[axe] (-0.18,0) -- (2.15,0) node[anchor=north east, inner sep=3pt, font=\small] {$t$};
  \draw[axe] (0,#3) -- (0,1.62);
  \node[anchor=south west, inner sep=1pt, font=\small] at (-0.1,1.62) {#2};
  #4
\end{scope}}

\node[font=\small\bfseries] at (3.6,2.35) {Uniform\'ement frein\'e};

% (1) position : croissante, concavité vers le bas, qui s'aplatit
\petitrepere{0}{$x$ ou $\theta$}{-0.18}{%
  \draw[courbe] plot[domain=0:1.95, samples=40, smooth] (\x,{1.42*\x-0.364*\x*\x});}

% (2) vitesse : droite décroissante qui coupe l'axe des temps
\petitrepere{2.7}{$v$ ou $\theta'$}{-0.95}{%
  \draw[courbe] (0,1.30) -- (1.95,-0.55);
  \fill[solideA] (1.37,0) circle (1.6pt);}

% (3) accélération : constante négative (droite sous l'axe des temps)
\petitrepere{5.4}{$a$ ou $\theta''$}{-0.95}{%
  \draw[courbe] (0,-0.62) -- (1.95,-0.62);
  \node[font=\small, text=solideA, anchor=north] at (1.0,-0.72) {$a < 0$};}
\end{tikzpicture}
\end{document}
